\documentclass[11pt]{article}
\usepackage[margin=1in]{geometry}
\usepackage{amsmath,amssymb,amsthm,hyperref}
\newtheorem{proposition}{Proposition}
\newcommand{\T}{\mathbb T}
\newcommand{\U}{\mathcal U}
\newcommand{\E}{\mathbb E}
\title{HOCF--RIESZ: Round 002\\Smooth iid residuals and quantitative interfaces}
\author{Campaign research memorandum}
\date{17 September 2026}
\begin{document}
\maketitle
\noindent Published input checkpoint:
\texttt{a06178658d1e3d458536ff312ca793947212ec67}.
This memorandum records bounded mathematical advances within the fixed smooth
iid model. The accompanying Round 002 report records the independent audit
verdicts and immutable proof hashes. The singular subcritical and critical
missions remain open; no cutoff removal is implicit in a fixed-smooth bound.

\section{Model and statement}
On the unit torus of volume one let $K=-\nabla g$ for a smooth real even
interaction $g$, and let $b=-\nabla V$ be smooth. For $N\ge2$, $\beta>0$,
\[
 dx_i=\left(b(x_i)+N^{-1}\sum_{j\ne i}K(x_i-x_j)\right)dt
             +\sqrt{2/\beta}\,dW_i.
\]
Initial positions are iid with smooth positive probability density $\mu_0$,
independently of the Brownian motions. The deterministic reference solves
\[
 \partial_t\mu=-\operatorname{div}((b+K*\mu)\mu)+\beta^{-1}\Delta\mu,
 \qquad \mu|_{t=0}=\mu_0.
\]
Write $\eta=N^{-1}\sum_i\delta_{x_i}$ and $\rho=\eta-\mu$.
The centered deleted-label statistic $\U_k$ alternates over all subsets of
$k$ slots, filling that subset by distinct particle labels with denominator
$N^{|A|}$ and the other slots by $\mu$. All tuples are ordered.
Set $P=\U_2/2$ and $\sigma=\sqrt N\min(\sqrt\beta,1)$.

Let $f$ be the full backward linearized one-body test and $\Phi$ its symmetric
zero-terminal pair corrector, with the operators frozen in Round 001. For
$r=4d+12$ suppose their $C^r$ norms and that of $\mu$ are bounded uniformly in
$N,\beta,t\in[0,T]$. Put
\[
 F=C\Phi=\operatorname{Sym}_3
       [K(x-z)\cdot\nabla_x\Phi(x,y)],
\]
where symmetrization averages six permutations. THM-010 establishes
\begin{align*}
 \E|P[\Phi_0]|&\le I/N,&
 \E\int_0^T|\U_3[F_t]|\,dt&\le TC_3/N^{3/2},\\
 \E[M_\Phi]_T&\le 2TH^2/(\beta N^2),&
 \E\operatorname{TV}[M_f,M_\Phi]&\le2TL_fH/(\beta N^{3/2}).
\end{align*}
All constants are explicit in the constructor memorandum, equations
(2.1)--(2.4), and are independent of $N,\beta$. They depend on the fixed
kernel, drift, horizon and displayed test norms. In particular the four
scaled targets have bounds, respectively,
\[
 O(N^{-1/2}),\qquad O(N^{-1}),\qquad O(N^{-1}),\qquad O(N^{-1/2}),
\]
when the first two are multiplied by $\sigma$ and the brackets by $\sigma^2$.
These are absolute estimates under the actual interacting law.

\section{Coupling, moments, and exact deletion}
Here is the load-bearing mechanism. Let $q=\lfloor d/2\rfloor+2$ and use
Fourier characters $e^{2\pi i k\cdot x}$. Define
\[
 \|\zeta\|_{-q}^2=\sum_{k\in\mathbb Z^d}
       (1+4\pi^2|k|^2)^{-q}|\widehat\zeta(k)|^2.
\]
The finite constants
\[
 B^2=\sum_k(1+4\pi^2|k|^2)^{-q},\qquad
 D^2=\sum_k4\pi^2|k|^2(1+4\pi^2|k|^2)^{-q}
\]
give $\|\delta_x\|_{-q}=B$ and
$\|\delta_x-\delta_y\|_{-q}\le D|\widetilde x-\widetilde y|$.
For independent samples of any probability law, expansion of the sixth
Hilbert moment shows
\[
 \left\|N^{-1}\sum_i\delta_{y_i}-\mu\right\|_{L^p(\Omega;H^{-q})}
       \le 2\,3^{7/6}B/\sqrt N,\qquad 1\le p\le6.
\]
Indeed a singleton index has zero expectation in the six-factor expansion;
a surviving list uses at most three labels. The second moment has the sharper
bound $B^2/N$.

Couple the particles to independent nonlinear particles driven by the same
initial samples and Brownian motions, with deterministic drift $b+K*\mu_t$.
For lifted displacements and their empirical mean $D_N$, the noises cancel.
Writing $\overline\rho$ for the independent empirical discrepancy gives
\[
 D_N(t)\le\mathcal K\int_0^t e^{L(t-s)}\|\overline\rho_s\|_{-q}\,ds,
 \qquad \|\rho_t\|_{-q}\le\|\overline\rho_t\|_{-q}+DD_N(t),
\]
where
$\mathcal K=(\sum_a\|K_a\|_{H^q}^2)^{1/2}$ and
$L=d(\|b\|_{C^1}+2\|K\|_{C^1})$.
Minkowski's inequality yields
\[
 \sup_{t\le T}\|\rho_t\|_{L^p(\Omega;H^{-q})}
 \le\frac{R}{\sqrt N},\qquad
 R=2\,3^{7/6}B\left(1+D\mathcal K\frac{e^{LT}-1}{L}\right),
\]
with the fraction interpreted as $T$ at $L=0$. The supremum is outside the
expectation. There is no claim of temperature-uniform supremum-in-time
empirical moments.

For symmetric kernels the exact, configuration-wise deletion formulas are
\begin{align*}
 \U_2[\Phi]&=\rho^{\otimes2}[\Phi]-N^{-1}\eta[\Phi(x,x)],\\
 \U_3[F]&=\rho^{\otimes3}[F]
 -3N^{-1}(\eta\otimes\rho)[F(x,x,y)]
 +2N^{-2}\eta[F(x,x,x)].
\end{align*}
The coefficients follow by inclusion-exclusion of label equalities, including
$N=2$. Smooth tensor Sobolev duality bounds the first cubic term by
$Q^3\|F\|_{C^{3q}}\|\rho\|_{-q}^3$, with $Q=(d+1)^{q/2}$.
The partial diagonal is bounded by $Q\|F\|_{C^q}\|\rho\|_{-q}$.
For the actual source $C\Phi$, the full diagonal is zero since $K(0)=0$.
These observations prove the cubic absolute bound without assuming the evolved
law is iid. The pair bound uses the initial second moment and retains its
nonzero deletion bias.

The exact pair gradient is $\nabla_{x_i}P=N^{-1}H_i$, where
\[
 H_i=\int\nabla_1\Phi(x_i,y)\rho(dy)-N^{-1}\nabla_1\Phi(x_i,x_i).
\]
Sobolev duality in $y$, uniformly in $x_i$, gives
$|H_i|\le L_\Phi\|\rho\|_{-q}+D_\Phi/N$. Hence
\[
 [M_\Phi]_T=\frac2{\beta N^2}\sum_i\int_0^T|H_i|^2dt,
 \qquad
 [M_f,M_\Phi]_T=\frac2{\beta N^2}\sum_i\int_0^T\nabla f(x_i)\cdot H_i\,dt.
\]
These yield the two bracket estimates. Their scaling uses exactly
$\sigma^2/\beta=N\min(1,\beta^{-1})$. No signed covariance cancellation is
needed. The two lower drift contractions have absolute integral at most
$5Td\|K\|_\infty\sup_t\|\Phi_t\|_{C^1}/N$.

\section{Qualification of the smooth norms}
The separate root lemma records bounds for existing smooth solutions in the
frozen model. Write $b_m=\|b\|_{C^m}$, $\kappa_m=\|K\|_{C^m}$ and
$U_m=b_m+\kappa_m$. Differentiating the forward density equation and applying
the finite-family maximum principle gives
\[
 \sup_{t\le T}\|\mu_t\|_{C^m}
 \le\|\mu_0\|_{C^m}
       \exp\{d(2^{m+1}-1)U_{m+1}T\}.
\]
The differentiated transport contributes $d(2^m-1)U_{m+1}$ and the differentiated
reaction contributes $d2^mU_{m+1}$. Nonnegative diffusion introduces no adverse
constant. For the one-body backward equation, derivatives in the response
operator fall on $K$; for $m\ge1$ its $C^m$ norm is at most $d\kappa_m$.
Thus, with $a_m=d(2^m-1)U_m+d\kappa_m$,
\[
 \sup_{t\le T}\|f_t\|_{C^m}\le e^{a_mT}\|h\|_{C^m}.
\]
The audited pair estimate THM-009 then implies, for $m\ge2$,
\[
 \sup_{t\le T}\|\Phi_t\|_{C^m}
 \le T e^{\overline c_mT}d2^{m+1}\kappa_m
                         e^{a_{m+1}T}\|h\|_{C^{m+1}},
\]
where
\begin{align*}
 \overline M&=d\|\mu_0\|_{C^1}e^{3dU_2T},\\
 \overline c_m&=2d(2^m-1)(b_m+3\kappa_m/2)
                      +2(d\kappa_m+\kappa_0\overline M).
\end{align*}
Taking $m=r$ supplies the declared regularity bounds for fixed smooth data,
uniformly in $N,\beta$. These constants diverge with a singular heat cutoff;
this qualification does not discharge a singular estimate.

\section{Finite-dimensional Gaussian criterion}
Fix a finite list $(t_a,\phi_a)_{a=1}^m$. Let $f_a^N$ denote the full
one-body backward test on $[0,t_a]$, with its pair corrector on that interval.
Put $a_N=\sigma_N/\sqrt N$ and $c_N=\sigma_N^2/(N\beta_N)$. The deterministic
covariance matrices are
\begin{align*}
 I_N^{ab}&=a_N^2\operatorname{Cov}_{\mu_0}
                         (f_a^N(0),f_b^N(0)),\\
 D_N^{ab}&=2c_N\int_0^{t_a\wedge t_b}
        \mu_r^N(\nabla f_a^N(r)\cdot\nabla f_b^N(r))\,dr.
\end{align*}
THM-011 proves that entrywise convergence $I_N\to I$, $D_N\to D$ implies
\[
 (\sigma_N\rho_{t_a}^N(\phi_a))_{a=1}^m
       \Longrightarrow\mathcal N_m(0,I+D).
\]
The Gaussian may be degenerate. The full proof is the sealed memorandum
\nolinkurl{ROUND_002_SMOOTH_GAUSSIAN.md}; its independent review is recorded
separately in the round report.

The exact corrected identity decomposes the vector as $Z_N=S_N+L_N+R_N$,
where $S_N$ is the initial linearized statistic, $L_N$ the leading one-body
martingale, and $\|R_N\|_2\le C N^{-1/2}$. The independent Fourier residual
proof gives the necessary $L^2$ bounds. Its expectation also proves
$|\E\rho_{t_a}^N(\phi_a)|\le C/N$, so the asserted deterministic mean-field
centering is valid at scale $\sigma_N$.

For a fixed vector $\theta$, let $Q_N^\theta$ be the bracket of
$\theta\cdot L_N$. Actual-law empirical control gives
\[
 \E|Q_N^\theta(T)-\theta^TD_N\theta|\le C_\theta N^{-1/2},
 \qquad 0\le Q_N^\theta(t)\le B_\theta,
\]
with a deterministic bound. Thus
$E_N(t)=\exp(i\theta\cdot L_N(t)+Q_N^\theta(t)/2)$ is a true bounded complex
martingale, and $\E[E_N(T)\mid\mathcal F_0]=1$. For every bounded
initial-measurable $Y_N$,
\[
 \left|\E[Y_Ne^{i\theta\cdot L_N}]
       -e^{-\theta^TD_N\theta/2}\E Y_N\right|
 \le C_\theta\|Y_N\|_\infty N^{-1/2}.
\]
An elementary third-order Taylor bound for the triangular iid initial sum
then yields
\[
 \E e^{i\alpha\cdot S_N+i\theta\cdot L_N}
 =e^{-\alpha^TI_N\alpha/2-\theta^TD_N\theta/2}
        +O_{\alpha,\theta}(N^{-1/2}).
\]
Independent Gaussian smoothing and the uniform second moments give weak
convergence, including degenerate covariance. This proves independence in
the limit, without asserting finite-particle independence.

For general fixed smooth data the full-response energy identity is
\[
 \frac d{dt}\mu_t(f_t^2)
 =2\beta_N^{-1}\mu_t(|\nabla f_t|^2)-2\mu_t(f_tRf_t).
\]
The uniform spatial bounds imply
$\int_0^{t_a}\mu_r(|\nabla f_a|^2)dr\le C\beta_N$.
Consequently $I_N,D_N\to0$ if $\beta_N\to0$, and the observable converges
to zero in $L^2$ at this normalization. If $\beta_N\to\infty$, the leading
noise has variance $O(\beta_N^{-1})$; identifying the remaining initial
covariance requires convergence of the backward kernels or their covariance.

For free heat, Haar background and a real cosine mode of eigenvalue
$\omega_k=4\pi^2|k|^2$, the exact covariance at times $t,s$ is
\[
 \frac{\min(\beta_N,1)}2
        e^{-\omega_k|t-s|/\beta_N}.
\]
The initial and noise contributions sum to this formula at every finite
$N$. At time zero, alternating $\beta_N=1/2$ and $\beta_N=2$ gives covariance
limits $1/4$ and $1/2$. Hence arbitrary oscillating temperatures need not
have a single limiting law. No field/path tightness follows from this
finite-list theorem.

\section{Independent all-order reconstruction}
AUD-004 compares an isolated reconstruction with the frozen R1 formulas.
Both give the same upward, internal and lower drift terms, with the apparent
factor $k(k-1)$ in the new lower kernel becoming $k$ when that kernel is
defined as the sum of $k-1$ incident pair terms. Every shared-label bracket
also agrees. In a centered representation, a matching of $p$ labels across
the two blocks has one distinguished Brownian edge, coefficient
$2/(\beta_NN^p)$, and multiplicity
\[
 p\binom{k}{p}\binom{\ell}{p}p!
 =k\ell\binom{k-1}{p-1}\binom{\ell-1}{p-1}(p-1)!.
\]
Every subset of the $p$ shared quotient slots is then integrated once against
$\mu$; the remaining slots enter the centered statistic. This is an exact
factorial-to-centered inversion, valid also for $k>N$. Individual centered
terms need not be positive. The finite algebra now has both isolated
reconstruction and hostile review passes. Critical effective power counting
remains a separate gate.

\section{All-order estimates and a sufficient cutoff condition}
THM-014 repairs THM-012's missing symmetry qualification. Its partition formula,
with full proof and constants
in the power-counting memorandum. For a partition $\pi$ of $[k]$, write
$j(\pi)$ for its singleton count and
\[
 a(\pi)=k-|\pi|,\qquad c(\pi)=\prod_{B\in\pi}(|B|-1)!,\qquad
 \Delta(\pi)=\frac12\sum_{|B|\ge2}(|B|-2).
\]
If $\Phi_\pi$ identifies variables inside each block, then
\[
 \U_k[\Phi]=\sum_{\pi}(-1)^{a(\pi)}c(\pi)N^{-a(\pi)}
 \int\Phi_\pi\prod_{|B|=1}\rho(dx_B)\prod_{|B|\ge2}\eta(dx_B).
\]
It holds also when $k>N$. The proof selects injective assignments by the
signed cycle-partition weights, then sums occupied-singleton versus background
choices to obtain $\rho$ in each remaining singleton slot.

Symmetrization with independent copies and signs gives the iid Hilbert bound
\[
 \E\|\overline\rho_t\|_{-q}^{2m}
 \le (2B)^{2m}(2m-1)!!N^{-m}.
\]
The same coupling transfers every fixed moment to the interacting law. With
$S=1+D\mathcal K(e^{LT}-1)/L$ and
$R_p=2BS\sqrt{2\lceil p/2\rceil}$, it gives
$\sup_t\|\|\rho_t\|_{-q}\|_{L^p}\le R_pN^{-1/2}$.
Tensor duality in singleton variables therefore proves
\[
 \|\U_k[\Phi_t]\|_{L^p}
 \le\sum_\pi c(\pi)(QR_{pj(\pi)})^{j(\pi)}
       \|\Phi_t\|_{C^{j(\pi)q}}
           N^{-k/2-\Delta(\pi)}.
\]
At $j=0$ the moment factor is one. Pair blocks have $\Delta=0$; larger
blocks have a strict extra power. The summed constant $C_{k,p}(N)$ is a finite
explicit partition sum, and obeys the conservative fixed-data bound
\[
 C_{k,p}(N)\le k!\exp\{2QBS\sqrt{pk+2}\}.
\]
The exact finite sum, rather than this last exponential majorant, is used
when the data norms vary with cutoff.

For rooted derivatives use symmetric kernels, replacing any kernel by its
averaged symmetrization; this preserves $\U_k$ and contracts every $C^m$ norm.
Excluding the root replaces $R_p$ by $R_p+B$ and gives constants
$D_{k-1,2}(N)$ defined by the same partition sum. In particular, if
$B_k'=\sup_t\|\Phi_{k,t}\|_{C^{(k-1)q+1}}$,
\[
 \E\operatorname{TV}[M_k,M_\ell]
 \le\frac{2Td k\ell}{\beta_N}
 B_k'B_\ell'D_{k-1,2}(N)D_{\ell-1,2}(N)
                  N^{-(k+\ell)/2}.
\]
The scaled power is $N^{1-(k+\ell)/2}\min(1,\beta_N^{-1})$.
The detailed generator table also retains the derivative cost of every
upward, internal, single-lowering and double-lowering term. A double
contraction is of the same raw order as $\U_k$; it is not a new gain at its
own level. Infinite summation requires the displayed order-dependent norms
and coefficients to be summable, including all drift and cross-bracket terms.
A separately stated lemma with geometric kernel-norm growth and coefficients
$1/k!$ proves statistical and martingale tail control; those extra hypotheses
are not asserted for the actual hierarchy.

For the frozen positive-power Riesz heat sequence, $0<s<d$, the report proves
explicit constants with
\[
 \|K_\varepsilon\|_{C^m}
 \le L_m\varepsilon^{-(s+m+1)/2},\qquad
 \|K_\varepsilon\|_{C^1}\ge c\varepsilon^{-(s+2)/2}.
\]
The gamma factors in $L_m$ and the lattice lower bound are displayed in the
proof. Thus the particular coupling constant grows exponentially in
$\varepsilon^{-(s+2)/2}$ on a fixed positive horizon.
For the regularized pair problem, a sufficient joint-limit condition is
\[
 \varepsilon_N^{-(s+2)/2}
 +\log\left(1+
   \sup_t\|\Phi_t^{N,\varepsilon_N}\|_{C^{3q+1}}+
   \sup_t\|f_t^{N,\varepsilon_N}\|_{C^1}\right)=o(\log N).
\]
All four regularized residual targets then vanish, uniformly over positive
temperatures. If those norm envelopes grow polynomially, a sufficient choice
is $\varepsilon_N=(\log N)^{-\gamma}$ with $0<\gamma<2/(s+2)$.
More generally finite uniform envelopes at each fixed cutoff permit a
sufficiently slow diagonal choice. The fixed-data qualification supplies such
finite envelopes for the existing smooth family; it does not imply polynomial
cutoff growth.

The missing comparison is a different assertion, for example
\[
 \sigma_N\left\langle\phi,
  (\eta_T^{\mathrm{sing}}-\mu_T^{\mathrm{sing}})
  -(\eta_T^{\varepsilon_N}-\mu_T^{\varepsilon_N})\right\rangle
          \longrightarrow0
\]
in a specified mode, with singular well-posedness justified. A cutoff slow
enough for these derivative estimates may smooth the microscopic scales
needed for that comparison. No comparison or critical summation is claimed.

\section{Diffusivity continuity and identified smooth limits}
The separate THM-013 proof constructs the smooth inviscid reference by the
characteristic fixed point
\[
 F_t(x)=x+\int_0^t\left[b(F_s(x))+
       \int K(F_s(x)-F_s(y))\mu_0(y)\,dy\right]ds.
\]
On periodic displacements its right side is Lipschitz with constant
$L_b+2L_K$. Bounded speed allows continuation on every finite horizon.
The resulting smooth flow has positive Jacobian determinant and pushes
$\mu_0$ to the unique smooth inviscid solution. Smooth positive-diffusivity
mean-field existence remains the declared model input. The full inviscid
backward equation is constructed by a Volterra series for transport plus
the bounded response; its $n$th term has the displayed factorial denominator.

For every fixed $m$ the proof establishes
\begin{align*}
 \sup_{t\le T}\|\mu_t^\nu-\mu_t^{\nu'}\|_{C^m}
       &\le C_{\mu,m}|\nu-\nu'|,\\
 \max_a\sup_{t\le t_a}\|f_a^\nu(t)-f_a^{\nu'}(t)\|_{C^m}
       &\le C_{f,m}|\nu-\nu'|,
\end{align*}
with explicit fixed-data constants and no inverse diffusivity, including
$\nu=0$. Two additional derivatives of the density or backward test are
required by the Laplacian forcing. In the backward difference, both the
velocity difference and
\[
 (R_{\mu^\nu}-R_{\mu^{\nu'}})f(y)
     =\int K(z-y)\cdot\nabla f(z)
                    (\mu^\nu-\mu^{\nu'})(z)\,dz
\]
are retained. Output derivatives hit $K$, so this last term needs only the
supremum norm of the density difference. Maximum-principle commutators and
Gronwall prove the stated bounds; there is no unproved parabolic smoothing
estimate at $\nu=0$.

It follows that the covariance matrices in THM-011 converge for every finite
positive limiting inverse temperature. Their limits are exactly the formulas
already displayed, evaluated at the limiting reference and full backward
kernels. If $\beta_N\to\infty$, then $\nu_N\to0$, and
\[
 I_N^{ab}\longrightarrow
       \operatorname{Cov}_{\mu_0}(f_a^0(0),f_b^0(0)),
 \qquad D_N\longrightarrow0.
\]
Combined with the high-temperature zero limit, these results identify every
physical-temperature convergent subsequence in the fixed-smooth finite-list
model. For $0<s<d$, a critical-labelled sequence
$\beta_N\sim\lambda N^{1-s/d}$ tends to infinity. This gives the inviscid
initial Gaussian for a \emph{fixed smooth} interaction only; it is not the
singular critical law.

Exact checks include free heat, zero/equal/different terminal times, both
signs of a one-mode interaction and its removable resonance. The full
one-mode response yields
\[ f_a^\nu(t,x)=e^{-(2\pi)^2(\nu+a/2)(t_a-t)}\cos(2\pi x). \]
A high-frequency free test also shows why a uniform Lipschitz estimate cannot
depend solely on the terminal $C^m$ norm without extra derivatives.

\section{Law-class obstruction}
For the smooth circle interaction $g(x)=a\cos(2\pi x)$ with $a<0$, canonical
Gibbs density is proportional to
$\exp(\beta |a|N|\eta(e_1)|^2/2)$. Translation invariance makes the one-point
marginal uniform. If every particle lies in a sufficiently short fixed arc
of length $p_\delta>0$, then $|\eta(e_1)|^2\ge1-\delta/2$. Its contribution
to the partition function gives
\[
 \mathbb P_{\rm Gibbs}(|\eta(e_1)|^2\le1-\delta)
 \le p_\delta^{-N}e^{-\beta|a|N\delta/4}.
\]
For every $\beta_N\to\infty$, the pair mode
$\U_2[\cos(2\pi(x-y))]=|\eta(e_1)|^2-1/N$ tends to one in probability.
Thus the iid moment estimates cannot be transferred to arbitrary smooth
Gibbs preparation, even with exact first-marginal centering. This hostile-
reviewed obstruction is outside the iid theorem and is not a counterexample
to a positive-Riesz singular statement.





\section{Scope and next load-bearing interface}
The old condition $\beta_NN^{2s/d-1}\to0$, full microscopic subcriticality
$\lambda_N=\beta_NN^{s/d-1}\to0$, and criticality
$\lambda_N\to\lambda\in(0,\infty)$ remain distinct. No Riesz exponent occurs
in the fixed-smooth coupling argument. Its higher-order and cutoff constants
must be controlled before they can decide singular critical survival,
truncation, or resummation. In particular an exponential in
$T\|\nabla K_\varepsilon\|_\infty$ is not an effective power of $\lambda_N$.

No assertion about Gibbs preparation, the physical gap process, exact-marginal
centering, source covariance normalization, or field/path tightness is added
by this memorandum. Those interfaces retain their separate campaign records.
\end{document}
